%! Unit 6 — Trigonometric Functions — Unit Cover Page
\documentclass[10pt]{article}
\usepackage{precalculus-article}
\usepackage{precalculus-boxes}
\usepackage{ltablex}
\keepXColumns

\begin{document}
\thispagestyle{empty}

% Full-bleed navy banner
\begin{tikzpicture}[remember picture, overlay]
  \fill[navy] (current page.north west)
    rectangle ([yshift=-1.15in]current page.north east);
\end{tikzpicture}
\vspace{-1.05in}

\begin{center}
  {\color{white}\textbf{\LARGE Precalculus}}\\[6pt]
  {\color{skymid}\Large\textbf{Unit 6: Trigonometric Functions}}\\[4pt]
  {\color{white}\small 8 Lessons + Unit Test \quad|\quad Class Periods: 18--22}
\end{center}

\vspace{0.26in}

\begin{tcolorbox}[colback=goldbg, colframe=goldacc, boxrule=0.8pt, arc=2mm,
    left=3mm, right=3mm, top=2mm, bottom=2mm]
\small\textbf{Unit Essential Questions:}
\begin{itemize}[leftmargin=1.4em, itemsep=2pt, topsep=2pt]
  \item How do sine and cosine turn motion around a circle into functions we can graph and analyze?
  \item How do we model periodic phenomena with sinusoidal functions?
\end{itemize}
\end{tcolorbox}

\vspace{0.14in}

\begin{tocbox}
\small
\renewcommand{\arraystretch}{1.45}
\begin{tabularx}{\linewidth}{@{} c >{\bfseries}X l @{}}
\rowcolor{sky}
\textbf{\#} & \textbf{Lesson Title} & \textbf{Content Source} \\
\midrule
6.1 & Periodic Phenomena & CED 3.1 \\
6.2 & Angles and Radian Measure & Split of CED 3.2 \\
6.3 & Sine, Cosine, and Tangent & CED 3.2 \\
6.4 & Sine and Cosine Function Values & CED 3.3 \\
6.5 & Sine and Cosine Function Graphs & CED 3.4 \\
6.6 & Sinusoidal Functions & CED 3.5 \\
6.7 & Sinusoidal Function Transformations & CED 3.6 \\
6.8 & Sinusoidal Function Context and Data Modeling & CED 3.7 \\
\midrule
\rowcolor{sky}
     & \textbf{Sample Test} &      \\
\end{tabularx}
\end{tocbox}

\vspace{0.14in}

\begin{skillbox}[Mathematical Practices --- Unit 6 Focus]{goldbox}
\renewcommand{\arraystretch}{1.4}
\begin{tabularx}{\linewidth}{@{} >{\bfseries\color{navy}}p{0.12\linewidth} X @{}}
MP 2.A & Identify information from graphical, numerical, analytical, and verbal representations. \\
MP 2.B & Construct equivalent graphical, numerical, analytical, and verbal representations. \\
MP 1.C & Construct new functions using transformations, compositions, inverses, or regressions. \\
MP 3.A & Describe the characteristics of a function with varying levels of precision. \\
\end{tabularx}
\end{skillbox}

\end{document}
